\uuid{UL3Y}
\titre{Nombres premiers}
\auteur{bodin ; croizat ; sacre}
\organisation{exo7}
\licence{CC BY-NC-SA 4.0}
\datecreate{2022-03-08}
\niveau{L1}
\module{Algèbre}
\chapitre{Arithmétique dans Z}
\sousChapitre{Nombres premiers, nombres premiers entre eux}
\difficulte{2}
\theme{Nombres premiers}
\exo7qcmid{hilisit.arithmetique.031}

\contenu{
\begin{question}
\qclassification{103.04}
    Soient $p$ un nombre premier et $a,b$ des entiers avec $p | ab$.
    Par application du lemme d'Euclide, quelles sont les affirmations vraies ?
    \begin{answers}
        \bad{$p$ divise $a$ et $p$ divise $b$.}
        \good{$p$ divise $a$ ou $p$ divise $b$.}
        \bad{$p$ divise $a$ ou $p$ divise $b$, mais pas les deux en même temps.}
        \bad{$p$ ne divise ni $a$, ni $b$.}
    \end{answers}
    \begin{explanations}
        Lemme d'Euclide : Si $p$ premier et $p | ab$, alors $p|a$ ou $p|b$.
        Les autres affirmations sont fausses. Voici des contre-exemples :
        $2 | (3 \times 4)$ mais $2$ ne divise pas $3$, mais divise bien $4$ ;
        $2 | (4 \times 6)$ et $2|4$ et $2|6$.
    \end{explanations}
\end{question}
}
